\section{DBの終了及び精算}

\subsection{DBの終了}

\begin{fillblank}[数理人試験 H30 問題1 設問1 A〜E]
  \noindent \textbf{◯DB法}
  \begin{lawstructure}
    \begin{lawarticle}{第八十三条}{確定給付企業年金の終了}
      規約型企業年金は、次の各号のいずれかに該当するに至った場合に終了する。

      \lawitem{一}
      次条第一項の規定による終了の承認があったとき。

      \lawitem{二}
      第八十六条の規定により規約の承認の効力が失われたとき。

      \lawitem{三}
      第百二条第三項又は第六項の規定により規約の承認が取り消されたとき。

      \lawparagraph{2}
      基金は、次の各号のいずれかに該当するに至った場合に解散する。この場合において、当該基金型企業年金は、終了したものとする。

      \lawitem{一}
      第八十五条第一項の認可があったとき。

      \lawitem{二}
      第百二条第六項の規定による基金の解散の（　　A　　）があったとき。
    \end{lawarticle}


    \begin{lawarticle}{第八十四条}{厚生労働大臣の承認による終了}
      事業主は、実施事業所に使用される（　　B　　）の（　　C　　）で組織する労働組合があるときは当該労働組合、
      当該（　　B　　）の（　　C　　）で組織する労働組合がないときは当該（　　B　　）の（　　C　　）を代表する者の同意を得たときは、厚生労働大臣の承認を受けて、規約型企業年金を終了することができる。

      \lawparagraph{2}
      前項の場合において、実施事業所が二以上であるときは、同項の同意は、各実施事業所について得なければならない。

      \lawparagraph{2}
      第五条第二項及び第三項の規定は、第一項の終了の承認があった場合について準用する。この場合において、同条第三項中「承認を受けた規約」とあるのは、「承認を受けた旨」と読み替えるものとする。
    \end{lawarticle}


    \begin{lawarticle}{第八十五条}{基金の解散}
      基金は、代議員会において（　　D　　）の（　　E　　）により議決したとき、又は基金の事業の継続が不可能となったときは、厚生労働大臣の認可を受けて、解散することができる。

      \lawparagraph{2}
      第五条第二項及び第三項の規定は、前項の解散の認可があった場合について準用する。この場合において、同条第三項中「承認を受けた規約」とあるのは、「認可を受けた旨」と読み替えるものとする。
    \end{lawarticle}
  \end{lawstructure}
\end{fillblank}

\begin{answerbox}
  \begin{enumerate}[A: ]
    \item 命令
    \item 厚生年金保険の被保険者
    \item 過半数
    \item 代議員の定数
    \item 四分の三以上の多数
  \end{enumerate}
\end{answerbox}


\begin{truefalse}[年金1 2018 問題1(1)\textcircled{5}]
  下線部の正誤を判定し、誤っている場合は正しい内容に改めなさい。
  \\

  規約型企業年金の終了のケースとして、法人の消滅等により規約の承認の効力が失われた
ときは、その事由（法人の消滅等）が発生してから\underline{60日以内}にその旨を厚生労働大臣に届け
出なければならない
\end{truefalse}

\begin{answerbox}
  \textbf{誤}

  30日以内
\end{answerbox}

\begin{truefalse}[年金2 2020 問題1(1)\textcircled{1}]
  下線部の正誤を判定し、誤っている場合は正しい内容に改めなさい。
  \\

  確定給付企業年金を実施する企業年金基金は、代議員会において\underline{出席した代議員の三分の
二以上}の多数により議決したとき、または基金の事業の継続が不可能となったときは、厚生
労働大臣の認可を受けて、解散することができる。
\end{truefalse}


\begin{answerbox}
  \textbf{誤}

  代議員の定数の四分の三以上
\end{answerbox}

\newpage

\subsection{終了時の掛金の一括拠出}

\begin{fillblank}[数理人試験 H30 問題1 設問1 F]
  \noindent \textbf{◯DB法}
  \begin{lawstructure}
    \begin{lawarticle}{第八十七条}{終了時の掛金の一括拠出}
      第八十三条の規定により確定給付企業年金が終了する場合において、当該終了する日における積立金の額が、当該終了する日を第六十条第三項に規定する事業年度の末日とみなして同項の規定に基づき算定した（　　A　　）を下回るときは、第五十五条第一項の規定にかかわらず、事業主は、当該下回る額を、掛金として一括して拠出しなければならない。
    \end{lawarticle}
  \end{lawstructure}
\end{fillblank}


\begin{answerbox}
  \begin{enumerate}[A: ]
    \item 最低積立基準額
  \end{enumerate}
\end{answerbox}


\newpage

\subsection{支給義務の消滅、清算}

出題例なし

\newpage


\subsection{残余財産の分配}

\begin{truefalse}[年金1 2018 問題2(1)\textcircled{4}]
  下線部の正誤を判定し、誤っている場合は正しい内容に改めなさい。
  \\

  終了した確定給付企業年金において、残余財産の額が終了日の最低積立基準額を上回る場
合、残余財産の額から終了日の最低積立基準額を控除した額の分配方法は、規約変更の同意
を得ることで、任意の方法に変更することが認められる。
\end{truefalse}

\begin{answerbox}
  \textbf{誤}

  分配方法は、加入者等に係る責任準備金の額又は最低積立基準額を勘案して公平かつ合理的に行われるものであること。
\end{answerbox}

\newpage